\section{总结与延伸}

\subsection{一张表回收全讲}

本讲从 poker search 出发，在 Diplomacy 中落成完整 agent。下面的表把每个模块要解决的问题、主要机制和失败模式放到一起，便于后续阅读多智能体系统或 reasoning 论文时复用。

\begin{longtable}{L{0.20\textwidth}L{0.30\textwidth}L{0.40\textwidth}}
\toprule
模块 & 解决的问题 & 机制与主要风险 \\
\midrule
Base policy & 快速给出合理候选 & 覆盖广但在关键局面可能粗糙；不能代替 search \\
Inference-time search & 用额外计算改善当前决策 & 展开、评估和选择；可能放大错误 evaluator \\
Human anchor policy & 维持人类可理解的 convention & imitation 提供先验；也会继承平均人类弱点 \\
piKL & 在价值与 human-likeness 间权衡 & 用 KL penalty 限制策略漂移；$\lambda$ 决定折中 \\
Intent model & 把计划变成可控生成条件 & 近程 action 清晰可查；长期目标表达不足 \\
Dialogue model & 实现 grounded communication & 在 state、history 与 intent 条件下生成候选 \\
Action predictor & 预测他人收到消息后的行为 & 连接语言与策略；对分布外人类可能失准 \\
Value filter & 拒绝战略上有害的消息 & 比较 counterfactual expected value；忽略未建模后果 \\
Closed loop & 让沟通结果回到下一次规划 & 每个事件后重估 belief；复杂度和延迟上升 \\
\bottomrule
\end{longtable}

\subsection{最终结论}

第一，训练规模和推理时计算是两条不同能力轴。Poker 与 Go 说明，强 base model 仍可能因缺少 search 而远离可达性能。第二，语言 agent 的核心不是生成，而是语言、belief 与 action 的因果闭环。第三，人类数据既是行为先验，也是社会 convention 的载体，但不能替代价值优化。第四，系统若无法评估某类长期风险，就应诚实承认盲区并限制 action space，而不是用流畅输出掩盖 evaluator 缺失。

\begin{importantbox}{最值得带走的一句话}
把更多计算花在正确的决策上，把世界建模、策略选择和语言表达交给有清晰契约的模块，再让输出后果回到下一轮规划——这就是本讲从 poker search 到 Cicero 的共同设计原则。
\end{importantbox}

\subsection{拓展阅读}

以下资料均在课堂日期前公开，并直接构成本讲证据链。建议先读 Cicero 系统论文，再用 piKL 论文理解 planner，最后回看 poker 与 AlphaGo 的 search 传统。

\begin{itemize}
  \item \href{https://www.science.org/doi/full/10.1126/science.ade9097}{Bakhtin et al., Human-level play in the game of Diplomacy by combining language models with strategic reasoning}：Cicero 的系统架构、在线结果、对话与限制。
  \item \href{https://proceedings.mlr.press/v162/jacob22a.html}{Jacob et al., Modeling Strong and Human-Like Gameplay with KL-Regularized Search}：piKL 的正式目标和 human-policy modeling。
  \item \href{https://proceedings.neurips.cc/paper/2021/hash/95f2b84de5660ddf45c8a34933a2e66f-Abstract.html}{Anthony et al., Learning to Play Diplomacy with No Human Supervision}：No-Press Diplomacy、自博弈和多均衡问题。
  \item \href{https://www.science.org/doi/10.1126/science.aao1733}{Brown and Sandholm, Superhuman AI for heads-up no-limit poker}：Libratus 与安全子博弈求解。
  \item \href{https://www.science.org/doi/10.1126/science.aay2400}{Brown and Sandholm, Superhuman AI for multiplayer poker}：Pluribus 与多人扑克。
  \item \href{https://www.science.org/doi/10.1126/science.aar6404}{Silver et al., Mastering the game of Go without human knowledge}：AlphaGo Zero 与 search ablation。
  \item \href{http://www.incompleteideas.net/IncIdeas/BitterLesson.html}{Richard Sutton, The Bitter Lesson}：可扩展 learning 与 search 的历史论证。
\end{itemize}

\subsection{课后自检}

\begin{importantbox}{八个自检问题}
\begin{enumerate}
  \item Exploitability 与对固定对手的胜率有什么区别？
  \item 为什么 slide 2 不能证明 search 永远等价于固定倍数的 base-model scaling？
  \item Diplomacy 的 support 规则为什么让语言直接进入状态转移？
  \item Intent conditioning 为什么能同时提高 controllability 和 perplexity？
  \item 纯 self-play 与纯 imitation 各自怎样失败？
  \item piKL 中 $\lambda$ 太大或太小分别会发生什么？
  \item Value-based filter 为什么仍无法可靠评估长期撒谎的代价？
  \item 一个开放世界 agent 要迁移 Cicero 架构，还缺哪些 transition、reward 和权限边界？
\end{enumerate}
\end{importantbox}

\end{document}
